\section{Носитель двойки и законы сохранения}\label{sec:carrier}

Прежде чем спускаться по цепи звеньев двигателя, нужно выделить то, что евклидов шаг
\emph{сохраняет}, и то, что \emph{разделяет} две стороны пары. Настоящий раздел фиксирует и то и
другое: двойка есть общий делитель, удерживающий стороны врозь, и одновременно инвариант, который
переносит спуск. Всё здесь элементарно --- делимость и кольцевая арифметика на чистом ядре
Lean~4 и \code{mathlib}~\cite{lean4,mathlib}, без анализа, распределения и сита --- и всякое
утверждение ниже \green{}~машинно доказано при стандартных аксиомах, кроме единственного открытого
узла в конце.

Всюду мы работаем с \emph{центром} $m\ge 1$ евклидовой пары, чьи \emph{стороны} суть кандидаты
$6m-1$ и $6m+1$; всякое простое $p>3$ лежит в классе $\pm1\pmod 6$, так что пара с разностью $2$
выше $3$ центрируется в кратном шести. Условие $1\le m$ делает $6m-1$ честным предшественником $6m$
в $\N$ (усечённое вычитание), так что $6m+1=(6m-1)+2$ есть равенство натуральных чисел.

\subsection{Носитель двойки}

\begin{definition}[общий делитель сторон]\label{def:shared-divisor}
Натуральное $p$ есть \emph{общий делитель сторон} центра $m$, если одновременно $p\mid(6m-1)$ и
$p\mid(6m+1)$.
\end{definition}

Стороны отстоят ровно на два, $(6m+1)-(6m-1)=2$, и эта фиксированная разность управляет всей их
совместной делимостью.

\begin{theorem}[\code{twin\_sides\_shared\_dvd\_two}]\label{thm:twin_sides_shared_dvd_two}
\green{} Пусть $m\ge 1$ и $p\mid(6m+1)$, $p\mid(6m-1)$. Тогда $p\mid 2$.
\end{theorem}

\emph{Идея доказательства.} Перепишем б\'{о}льшую сторону через меньшую, $6m+1=(6m-1)+2$ (закрывается
\code{omega} над $\N$ при $m\ge1$). Тогда $p\mid\bigl((6m-1)+2\bigr)$, и так как $p\mid(6m-1)$,
свойство делимости суммы с делимым слагаемым (\code{Nat.dvd\_add\_right}) вынуждает $p\mid 2$.
Работает только \emph{разность} сторон: по существу это утверждение $\gcd(6m-1,\,6m+1)\mid 2$.

Так как у двойки лишь делители $1$ и $2$, ни одно простое выше двух не может занять обе стороны
сразу.

\begin{theorem}[\code{no\_large\_shared\_divisor}]\label{thm:no_large_shared_divisor}
\green{} Пусть $m\ge 1$ и $p>2$. Тогда $p\mid(6m+1)$ и $p\mid(6m-1)$ не могут выполняться
одновременно.
\end{theorem}

\emph{Идея доказательства.} По \cref{thm:twin_sides_shared_dvd_two} $p\mid 2$, откуда $p\le 2$ по
оценке делителя (\code{Nat.le\_of\_dvd}, положительность через \code{decide}); это противоречит
$2<p$ (\code{omega}). Это \emph{эксклюзивность}: простое $p>2$ делит не более одной стороны. Обе
теоремы \srcpath{Engine/Carrier.lean} доказаны в чистом ядре Lean~4, без \code{mathlib}, на трёх
примитивах. Эксклюзивность позволяет считать шесть активных множителей rank-$(3,3)$ разложения
различными и служит арифметическим источником отрицательной ассоциации в
\cref{thm:exclusive_no_backward}. Двойка --- это \emph{носитель}: она разделяет стороны наверху и
переносится вниз, сохраняющаяся величина, двойственная строго убывающей высоте из
\cref{sec:engine}.

\subsection{Первый закон: сохранение двойки}

При переходе к \emph{мультипликативным} координатам --- разложении сторон в произведения активных
простых слоя $(A,A^2]$ --- разделяющая двойка не размывается, а проступает как \emph{детерминант}
линейной системы, связывающей множители. Детерминант нельзя подвинуть: он равен $2$ (или
фиксированному кратному), либо решения нет. Получающееся семейство тождеств мы называем
\emph{первым законом}; во всех координатах оно имеет вид $XY-ZW=2K$, базовый случай $K=1$ даёт
$XY-ZW=2$. Термин ``закон'' употреблён в физическом смысле сохраняющегося заряда, а не логической
аксиомы.

\begin{definition}[rank-$(3,3)$ центр]\label{def:rank33}
Центр $m$ есть \emph{rank-$(3,3)$}, если существуют натуральные $a\le b\le v$ и $q\le r\le s$ с
\begin{equation}\label{eq:rank33}
6m-1=abv,\qquad 6m+1=qrs.
\end{equation}
\end{definition}

\begin{theorem}[\code{det\_law\_rank33}]\label{thm:det_law_rank33}
\green{} Пусть $1\le m$, $6m-1=abv$ и $6m+1=qrs$. Тогда $qrs=abv+2$.
\end{theorem}

\emph{Идея доказательства.} Разность $qrs-abv$ равна разности сторон $(6m+1)-(6m-1)=2$, но теперь
выражена через множители, а не через $m$; в Lean система линейно влечёт $qrs=abv+2$ (\code{omega},
с $1\le m$, снимающим неоднозначность усечённого вычитания). Это первая жёсткость: пять из шести
множителей $a,b,v,q,r,s$ определяют шестой. Мы фиксируем лишь жёсткость --- мы \emph{не} считаем
rank-$(3,3)$ центры и никогда не выдаём одно за другое.

Свёртка по одному множителю с каждой стороны в продукт-переменную превращает двойку в буквальный
детерминант $2\times2$.

\begin{definition}[продукт-координаты]\label{def:product-coords}
Для rank-$(3,3)$ центра положим $C=av$ и $D=qs$, так что стороны суть $Cb=6m-1$ и $Dr=6m+1$ со
свободными множителями $b,r$.
\end{definition}

\begin{theorem}[\code{det\_law\_CD}]\label{thm:det_law_CD}
\green{} Над $\Z$: если $Cb=6m-1$ и $Dr=6m+1$, то
\begin{equation}\label{eq:detCD}
Cb-Dr=-2.
\end{equation}
\end{theorem}

\emph{Идея доказательства.} Вычитание гипотез: $Cb-Dr=(6m-1)-(6m+1)=-2$ (\code{linear\_combination
h1 - h2}). Знак $-2$ --- ориентированный детерминант строк $(C,-D)$ против $(b,r)^\top$. При
$\gcd(C,D)=1$ \cref{eq:detCD} редуцируется по модулю $D$ к $b\equiv-2\,\overline{C}\pmod D$:
свободный множитель $b$ лежит в \emph{одном} классе вычетов по модулю $D=qs$. Это ``zero-one slot''
за верхними оценками счёта --- если $b$ загнан в интервал короче $D$, в фиксированном боксе
$(a,v,q,s)$ допустимо не более одного $b$. Мы фиксируем алгебраическую жёсткость; переход от неё к
счёту центров --- отдельный, аналитически тяжёлый шаг.

Первый закон --- про сохранение \emph{при движении}. Вдоль луча спуска двойка переносится с точки
на точку контролируемо.

\begin{definition}[луч и его шаг]\label{def:ray}
\emph{Лучом} назовём семейство точек с общими $v,s$, где $i$-я удовлетворяет rank-one уравнению
$u_i v+2=r_i s$, а переход задан целым шагом $u_2-u_1=sK$.
\end{definition}

\begin{theorem}[\code{det\_collapse}]\label{thm:det_collapse}
\green{} Пусть $u_1 v+2=r_1 s$, $u_2 v+2=r_2 s$, $u_2-u_1=sK$ и $s\ne 0$. Тогда
\begin{equation}\label{eq:detcollapse}
u_2 r_1-u_1 r_2=2K.
\end{equation}
\end{theorem}

\emph{Идея доказательства.} Домножим цель на $s$ и подставим оба rank-one уравнения:
$s(u_2 r_1-u_1 r_2)=u_2(u_1 v+2)-u_1(u_2 v+2)=2(u_2-u_1)=2sK$; слагаемые с $v$ сокращаются --- вот
где двойка выживает --- а \code{mul\_left\_cancel\char`_0\ hs} снимает $s\ne 0$. Название точно:
$u_i,r_i$ могут расти, но их перекрёстный детерминант схлопывается к строго линейной по $K$
функции с коэффициентом $2$. Спуск --- это шаг $K=1$, несущий ровно одну двойку.

Спуск по природе \emph{квадратичен}: евклидово деление сравнивает сторону с квадратом масштаба.
Следующее тождество --- диагональное ядро самоподобия.

\begin{definition}[евклидово разложение шага]\label{def:euclid-decomp}
При $abv+2=qP$ разложим $P=v^2+\Delta$ и $ab=qv+\Xi$ по масштабу $v$; остатки $\Delta,\Xi$ суть
\emph{евклидовы остатки} шага.
\end{definition}

\begin{theorem}[\code{euclid\_descent\_identity}]\label{thm:euclid_descent_identity}
\green{} При $abv+2=qP$, $P=v^2+\Delta$ и $ab=qv+\Xi$ выполнено
\begin{equation}\label{eq:euclid-descent}
q\,\Delta-v\,\Xi=2.
\end{equation}
\end{theorem}

\emph{Идея доказательства.} Раскрытие $qP=qv^2+q\Delta$ и $abv=qv^2+\Xi v$ сокращает $qv^2$,
оставляя $\Xi v+2=q\Delta$ (\code{linear\_combination -h1 - q*h2 + v*h3}). Это диагональное ядро
самоподобного спуска: после одного евклидова шага двойка воспроизводится на меньшей паре
$(\Delta,\Xi)$, так что спуск отображает ``двойку на $(a,b,v,q)$'' в тождественную меньшую задачу
``двойку на $(\Delta,\Xi)$''. Ни в одной форме правая двойка не обращается в ноль нетривиально ---
это противоречило бы $\gcd\mid 2$ и означало бы совпадение сторон; эта невозможность --
алгебраический предшественник необратимости.

Вырожденный, но полезный случай: один атом $a$ появляется дважды и две точки связаны через общий
мост $qs$. Здесь двойка переносится как масштабированная тень $12h=6\cdot 2h$, где шестёрка --- от
шага центровой решётки.

\begin{definition}[bridge-уравнение и его сдвиг]\label{def:bridge}
Повторный атом $a$ связывает две точки уравнениями $a B_1=qs\,Q_1-2$ и $a B_2=qs\,Q_2-2$ со сдвигом
$B_2-B_1=6\,(qs)\,h$.
\end{definition}

\begin{theorem}[\code{small\_det\_collapse}]\label{thm:small_det_collapse}
\green{} При $a\ne 0$, $qs\ne 0$ и указанных bridge-уравнениях
\begin{equation}\label{eq:smalldet}
B_2 Q_1-B_1 Q_2=12\,h.
\end{equation}
\end{theorem}

\emph{Идея доказательства.} Два сокращения. Сначала $qs(Q_2-Q_1)=qs\cdot 6ah$ даёт сопряжённый сдвиг
$Q_2-Q_1=6ah$ (сокращение на $qs\ne 0$). Затем
$a(B_2 Q_1-B_1 Q_2)=Q_1(aB_2)-Q_2(aB_1)+2(Q_2-Q_1)=a\cdot12h$, и сокращение на $a\ne 0$ даёт
\cref{eq:smalldet}. Двойка из моста и шестёрка из решётки перемножаются в $12$; условия ненулевости
--- содержательное требование \emph{активности} атома и моста. Собрав
\cref{thm:det_law_rank33,thm:det_law_CD,thm:det_collapse,thm:euclid_descent_identity,%
thm:small_det_collapse}, видим один скелет $XY-ZW=2K$: двойка не создаётся и не уничтожается
двигателем, а лишь переносится. Все доказаны без \code{sorry}; первый закон \emph{не} содержит
никакого счётного или распределительного утверждения.

\subsection{Спуск и граничный закон}

Пусть центр $m$ лежит в чистом ядре $\Om_A$ (обе стороны из простых $>A$) и одна сторона составна в
слое, неся активный множитель $a>A$:
\begin{equation}\label{eq:descent-step}
6m+\sigma=a\,(6n+\varepsilon),\qquad a>A,\ \varepsilon\in\{-1,+1\}.
\end{equation}
Удаляя $a$, спущенным центром объявляем то $n$ со стороной-остатком $6n+\varepsilon$. Так как $a>A$,
высота строго падает: $n<m/A$ --- поэлементная арифметика делимости, не оценка в среднем, и именно
она обрывает спуск ($m/A^t<1$ в конце концов). Введём $U:=6n+\varepsilon$.

\begin{theorem}[\code{two\_carry\_opposite}]\label{thm:two_carry_opposite}
\green{} Если $U=6n+\varepsilon$, то
\begin{equation}\label{eq:two-carry}
6n-\varepsilon=U-2\varepsilon.
\end{equation}
\end{theorem}

\emph{Идея доказательства.} $6n-\varepsilon=(6n+\varepsilon)-2\varepsilon=U-2\varepsilon$
(\code{linear\_combination -hU}). Зазор $2$ инвариантен при спуске: активный множитель $a$ унесён
из стороны-произведения, а разделяющая двойка перекочевала на противоположную сторону
$U-2\varepsilon$. То, что удерживает стороны врозь наверху ($\gcd\mid 2$), --- ровно то, что
переносится вниз.

Сторона-произведение $U=b\cdot v$ наследует множители $b,v>A$ из родословной $6m+\sigma$, так что
никакой старый малый простой её не портит.

\begin{theorem}[\code{no\_small\_divisor}]\label{thm:no_small_divisor}
\green{} Пусть $b,v$ --- простые с $A<b$, $A<v$, и $p$ --- простое с $p\le A$. Тогда
$p\nmid(b\cdot v)$.
\end{theorem}

\emph{Идея доказательства.} Если $p\mid b\cdot v$, то по лемме Евклида (\code{Nat.Prime.dvd\_mul})
$p\mid b$ или $p\mid v$; делимость простого на простое есть равенство
(\code{Nat.prime\_dvd\_prime\_iff\_eq}), значит $p=b$ или $p=v$, что противоречит $p\le A<b,v$
(\code{omega}). Так разграничена ответственность: если чистота нарушается неким $p\le A$, виновата
лишь противоположная сторона.

Называя обструкцией событие ``$p$ делит одну из сторон'' (boundary-exit из $\Om_A$), при $p\nmid U$
дизъюнкция схлопывается.

\begin{theorem}[\code{obstruction\_on\_opposite}]\label{thm:obstruction_on_opposite}
\green{} Пусть $\mathit{sideProd}=U$, $\mathit{sideOpp}=U-2\varepsilon$ и $p\nmid U$. Тогда
\begin{equation}\label{eq:obstruction}
\bigl(p\mid\mathit{sideProd}\ \lor\ p\mid\mathit{sideOpp}\bigr)\iff p\mid\mathit{sideOpp}.
\end{equation}
\end{theorem}

\emph{Идея доказательства.} Если $p\mid\mathit{sideProd}=U$, это противоречит $p\nmid U$; обратная
импликация тривиальна. В соединении с \cref{thm:no_small_divisor} (дающей $p\nmid U$ для малого
$p\le A$) и \cref{thm:two_carry_opposite} (опознающей $\mathit{sideOpp}=U-2\varepsilon$ как
перенесённую двойку) это даёт \emph{граничный закон}: единственная делимость, способная выбить
спущенный центр из $\Om_A$, есть $p\mid 6n-\varepsilon$. Вся комбинаторика обструкции сводится к
одной делимости на перенесённой двойке. Один такт двигателя таков: активный множитель уходит,
высота строго падает ($n<m/A$), двойка переносится (\cref{thm:two_carry_opposite}), а любой выход
опознаётся $p\mid 6n-\varepsilon$ (\cref{thm:obstruction_on_opposite}). Это закрывает
\emph{механику} одного такта; конечное дерево спусков, оканчивающееся граничными листьями, остаётся
комбинаторно возможным, и глобальный non-cover --- отдельный, более поздний узел.

\subsection{Нет хода назад на двух точках}

Рассматривая двигатель сразу на \emph{обеих} сторонах, эксклюзивность
(\cref{thm:no_large_shared_divisor}) запрещает ход назад и здесь, чисто арифметически.

\begin{definition}[индикаторы, ранги, эксклюзивность]\label{def:indicators}
Для простого $p$ с индексом $i$ из конечного $s$ положим $X,Y:\iota\to\N$,
$X_p=\mathbf{1}[\,p\mid 6m-1\,]$ и $Y_p=\mathbf{1}[\,p\mid 6m+1\,]$, с рангами
$r_-=\sum_{p\in s}X_p$ и $r_+=\sum_{p\in s}Y_p$. \emph{Эксклюзивность} (гипотеза \code{hexcl}) есть
$X_p\cdot Y_p=0$ для каждого $p\in s$.
\end{definition}

Утверждения ниже используют лишь $X_p Y_p=0$ и потому справедливы для любых неотрицательных весов с
исчезающим произведением; это абстрактная алгебра, не привязанная к индикаторам.

\begin{theorem}[\code{exclusive\_diag\_zero}]\label{thm:exclusive_diag_zero}
\green{} Для конечного $s$ и $X,Y:\iota\to\N$ с $X_p Y_p=0$ при всех $p\in s$,
\begin{equation}\label{eq:diag-zero}
\sum_{p\in s}X_p\cdot Y_p=0.
\end{equation}
\end{theorem}

\emph{Идея доказательства.} Сумма нулей (\code{Finset.sum\_eq\_zero hexcl}). Величина --- это
\emph{диагональ}: она считает простые, делящие обе стороны, --- self-переход ``в себя'', --- и
эксклюзивность стирает её поточечно, $\mathrm{DIAG}=0$.

\begin{theorem}[\code{exclusive\_no\_backward}]\label{thm:exclusive_no_backward}
\green{} Для конечного $s$ с $[\mathrm{DecidableEq}\ \iota]$ и $X,Y:\iota\to\N$ при эксклюзивности
\begin{equation}\label{eq:no-backward}
\Bigl(\sum_{i\in s}X_i\Bigr)\cdot\Bigl(\sum_{j\in s}Y_j\Bigr)=\sum_{i\in s}\sum_{j\in s\setminus\{i\}}X_i\cdot Y_j.
\end{equation}
\end{theorem}

\emph{Идея доказательства.} Раскроем произведение (\code{Finset.sum\_mul}, \code{Finset.mul\_sum})
в $\sum_i\sum_j X_i Y_j$, выделим член $j=i$ (\code{Finset.add\_sum\_erase}), равный $X_i Y_i=0$ по
эксклюзивности, а \code{zero\_add} оставляет внедиагональную часть. Произведение рангов
$r_-\cdot r_+$ целиком внедиагонально: self-переход $2\to2$ на обеих точках запрещён, остаются лишь
перекрёстные $i\ne j$. Раскладывая ковариацию
$\mathrm{Cov}(r_-,r_+)=\mathrm{CROSS}-\mathrm{DIAG}$, эксклюзивность обнуляет $\mathrm{DIAG}$
поточечно и смещает ассоциацию к отрицательной; численно ранг-$2$ поправка отрицательна на всех
проверенных масштабах. Но \cref{thm:exclusive_no_backward} доказывает лишь \emph{структуру} --- это
ещё не $\mathrm{Cov}(r_-,r_+)<0$ для реальных счётов, что требует $\mathrm{DIAG}>\mathrm{CROSS}$.
Это неравенство --- открытый узел.

\begin{conjecture}[отрицательная ассоциация рангов]\label{cnj:neg-assoc}
\red{} Ранги $(r_-,r_+)$ отрицательно ассоциированы на всех масштабах:
\begin{equation}\label{eq:neg-assoc}
N_{00}N_{33}\le N_{03}N_{30},
\end{equation}
эквивалентно $\mathrm{DIAG}\ge\mathrm{CROSS}$.
\end{conjecture}

Эксклюзивность даёт продуктовую CRT-структуру $\prod_p(c_p+a_p x+b_p y)$ без перекрёстного члена
$xy$, для которой отрицательная ассоциация элементарна; открытым остаётся остаток
модель$\to$реальность --- что реальный $\mathrm{CROSS}$ не превосходит $\mathrm{DIAG}$. Численно это
knife-edge ($1-R_{\mathrm{fc}}\to 0$) у стены чётности, и ни один имеющийся приём не даёт
распределённо-независимого закрытия --- поэтому узел объявлен открытым, а не закрытым.
